\chapter{The C interface: orders a robot can trust}

A unit in the field must not obey an order just because it arrived. The
command post signs every order under a key the unit was given at enlistment,
and the unit refuses anything that does not verify. The cryptography comes
from OpenSSL, and reaching it is the job of Poplog's foreign-function
interface --- which has existed since 1987 and is unusually deep: external code is loaded into the \emph{running} image by Poplog's own
dynamic linker, declared with a syntax form, and called with argument
conversion that understands Poplog's data representation. This chapter shows
the path this fork actually uses in production libraries, then names the
alternatives and the traps.

\section{The mechanisms, briefly}

\begin{center}
\footnotesize
\begin{tabular}{lp{8.2cm}}
\toprule
\code{external\_do\_load} & The primitive: a mark item, a list of object
  files or shared libraries, and a list of symbol specifications. Procedural.\\
\code{exload} & The declarative syntax form over it --- name the library,
  name the functions, give their types. \textbf{Start here.}\\
\code{exacc} & The syntax form that \emph{calls} an external function or
  accesses external data, converting arguments and results.\\
\code{cons\_access} & Builds an accessor procedure for a field of an external
  structure.\\
\code{lib external}, \code{lib newexternal} & Older higher-level layers; see
  \code{help external} and \code{help newexternal}.\\
\bottomrule
\end{tabular}
\end{center}

The reference material is \code{ref external} (mechanisms and data
representation), \code{ref defstruct} (structures) and \code{ref
external\_data}.

\section{The shim discipline}

The libraries in this fork --- \code{lib crypto} over OpenSSL, \code{popcurl}
over libcurl, \code{poppcre} over PCRE --- all go through a small C shim
rather than binding the library's own API directly. That is a
deliberate rule, and it is worth stating plainly:

\begin{callout}{Bind a shim, not the library}
Keep the \code{exload} surface to a handful of \textbf{non-variadic}
functions with \textbf{explicit lengths} and \textbf{caller-allocated
buffers}, and keep object lifecycle entirely on the C side. Variadic calls
through the FFI are hazardous --- notably on Apple \code{arm64}, where the
variadic calling convention differs from the regular one --- and library
handle lifetimes are exactly what a garbage collector should not be guessing
about.
\end{callout}

A shim is usually fifty lines. Here are two of the seven functions behind
\code{lib crypto}:

\begin{lstlisting}[language=C]
/* pop/extern/popcrypto/popcrypto_shim.c */
#include <openssl/evp.h>
#include <openssl/hmac.h>

/* digest of data[0..len) with a named algorithm; writes raw bytes to
   out, returns the digest length, or -1 */
int pcr_digest(const char *alg, const void *data, long len,
               unsigned char *out)
{
    const EVP_MD *md = EVP_get_digestbyname(alg);
    unsigned int n = 0;
    if (!md) return -1;
    if (!EVP_Digest(data, (size_t)len, out, &n, md, NULL)) return -1;
    return (int)n;
}

int pcr_random(unsigned char *buf, long n)
{
    return RAND_bytes(buf, (int)n) == 1 ? 0 : -1;
}
\end{lstlisting}

Note the shape: no allocation crosses the boundary, every buffer has an
explicit length, and every failure is a negative return rather than an
exception or an \code{errno} the caller must remember to read.

\section{Declaring it: \code{exload}}

\begin{pop}
uses crypto;   ;;; ...is, inside, this:

exload popcrypto
    #_IF sys_file_exists(sysfileok('$usepop/pop/extern/popcrypto/popcrypto.dylib'))
    ['$usepop/pop/extern/popcrypto/popcrypto.dylib']
    #_ELSE
    ['$usepop/pop/extern/popcrypto/popcrypto.so']
    #_ENDIF
(language C)
    lconstant
        pcr_digest(alg, data, len, out)             :int,
        pcr_hmac(alg, key, keylen, data, len, out)  :int,
        pcr_random(buf, n)                          :int,
        pcr_pbkdf2(alg, pass, passlen, salt, saltlen, iters, out, outlen) :int,
        pcr_version()                               :exptr,
        ;;; ...and pcr_gcm_encrypt / pcr_gcm_decrypt
    ;
endexload;
\end{pop}

Three things are happening. \code{popcrypto} is the \emph{mark item} --- the
tag under which this set of symbols is loaded and can be unloaded.
\code{\#\_IF} is a compile-time conditional, so one source file picks the
right shared-library extension on macOS and Linux. And each declaration gives
the argument names and the \emph{result} type: \code{:int} for a C \code{int},
\code{:exptr} for a pointer (Poplog wraps it in an external-pointer record;
\code{exacc\_ntstring} pulls a C string out of one). The full type vocabulary
--- \code{:byte}, \code{:short}, \code{:int}, \code{:long}, \code{:sfloat},
\code{:dfloat}, pointers, arrays and structures --- is in \code{ref external}.

\section{Calling it: \code{exacc}}

\begin{pop}
lconstant OUTMAX = 64;              ;;; SHA-512 is the widest digest

define crypto_digest(alg, data) -> raw;
    lvars n, out = inits(OUTMAX);   ;;; caller allocates the buffer
    unless isstring(alg) and isstring(data) then
        mishap(alg, data, 2, 'crypto_digest: two strings needed')
    endunless;
    exacc pcr_digest(cstr(alg), data, length(data), out) -> n;
    if n < 0 then
        mishap(alg, 1, 'crypto_digest: unknown algorithm')
    endif;
    substring(1, n, out) -> raw;
enddefine;
\end{pop}

\begin{repl}
: uses crypto;
: crypto_digest_hex('sha256', 'hello poplog') =>
** 8b05ad6e876a9ce82062d33918c82d6b9bfac5b0f1cb575bb4dc8cd70762c6ae
: crypto_version() =>
** OpenSSL 3.6.3 9 Jun 2026
\end{repl}

On top of that, the signed order is a dozen lines --- \ghfile{signed\_orders.p}:

\begin{pop}
vars fleet_key = 'a shared secret, provisioned at enlistment';

;;; wire format:  <hex mac> <space> <order text>
define sign_order(text) -> wire;
    crypto_hmac_hex('sha256', fleet_key, text) <> ' ' <> text -> wire
enddefine;

define verify_order(wire) -> text;
    lvars sp = locchar(` `, 1, wire), mac, body;
    false -> text;
    if sp then
        substring(1, sp - 1, wire) -> mac;
        allbutfirst(sp, wire) -> body;
        if mac = crypto_hmac_hex('sha256', fleet_key, body) then body -> text endif
    endif
enddefine;
\end{pop}
\begin{repl}
on the wire: 81bdc45a...581146c3 unit r1 advance to ridge
genuine:     unit r1 advance to ridge
tampered:    <false>          <- one character of the order changed
replayed:    <false>          <- a real signature on a different order
\end{repl}

\subsection{Two traps that bite immediately}

\textbf{Pop-11 strings are not null-terminated.} They are counted byte
vectors, so a string passed where C expects \code{char~*} must be terminated
explicitly. \code{lib crypto} does it in one line:

\begin{pop}
define lconstant cstr(s);
    s <> consstring(0, 1)       ;;; append a NUL
enddefine;
\end{pop}

Data arguments do not need this --- they are passed with an explicit length,
which is the better interface anyway.

\textbf{Packed fields versus full fields.} Poplog integers and floats are
tagged in ordinary records and vectors. Passed as direct arguments they are
converted for you; passed \emph{inside} a structure they are not, and the C
side sees encoded nonsense. An integer array destined for C must therefore be
built with a packed representation (\code{intvec} and friends), not
\code{newarray}. Strings made with \code{inits(n)} are packed byte vectors,
which is exactly why they serve as C buffers above.

\section{Building and shipping the shim}

The shim is built by a small script, not by the main \code{make}, so a Poplog
without OpenSSL still builds:

\begin{sh}
tools/build-popcrypto.sh      # -> pop/extern/popcrypto/popcrypto.{so,dylib}
\end{sh}
\begin{lstlisting}[language=bash]
if command -v pkg-config >/dev/null 2>&1 && pkg-config --exists libcrypto; then
    flags="$(pkg-config --cflags --libs libcrypto)"
else
    flags="-lcrypto"
fi
cc -O2 -shared -fPIC -o "$lib" "$src" $flags
\end{lstlisting}

And the library fails loudly at load time if the shim is missing, rather than
mysteriously at the first call:

\begin{pop}
unless sys_file_exists(sysfileok('$usepop/pop/extern/popcrypto/popcrypto.so'))
or    sys_file_exists(sysfileok('$usepop/pop/extern/popcrypto/popcrypto.dylib'))
then
    mishap(0, 'lib crypto: shim not built -- run tools/build-popcrypto.sh first')
endunless;
\end{pop}

Both \code{lib json} and its FFI-based sibling \code{lib crypto} are validated
on x86-64, AArch64 (Raspberry Pi 5) and RISC-V; \code{tools/test-crypto.sh}
is the acceptance suite and \code{tools/test-remote.sh \emph{host}} reruns it
anywhere with a built tree. That matters more than usual for FFI code,
because calling conventions are exactly what differs between ports --- the
RISC-V float ABI, for instance, was found and fixed by running these suites.

\section{Images, handles, and when not to bother}

External state does not survive \code{syssave}. Library handles, open
connections and \code{exptr} values are stale after a restore --- the right
shape for a binding is to \emph{mishap cleanly} on a stale handle rather than
crash, and the fix is always to reopen. Procedures, by contrast, survive
fine, so a restored image is ready to work as soon as you reconnect its
resources.

Finally: the FFI is not always the answer. Poplog can stream the output of a
shell command as a line repeater, and for a one-off that is often enough:

\begin{pop}
vars r = sys_obey_linerep('set +e; curl -s https://example.com/api | jq -r ".id"');
\end{pop}

Note \code{set~+e}: the command line runs under \code{errexit}, so without it
the first failing command aborts the whole pipeline \emph{silently} --- empty
output and no error.

The trade-off against the FFI is a process per call, and it is not subtle: on
the machine this chapter was written on, 100 round trips through
\code{sys\_obey\_linerep('echo hi')} took 284\,ms of wall time --- about
2.8\,ms each --- against an unmeasurable zero for 100 in-image calls. Use a
shim when the call sits in a loop; use the pipe when it does not.
